\documentclass[11pt]{article}
\usepackage[utf8]{inputenc}
\usepackage[margin=1in]{geometry}
\usepackage{amsmath}
\usepackage{graphicx}
\usepackage{booktabs}
\usepackage{hyperref}
\usepackage[round]{natbib}
\usepackage{float}
\usepackage{multirow}
\usepackage{array}

\title{Antipsychotic Drugs Selectively Decorrelate Long-Range Interactions in Deep Cortical Layers}

\author{%
  Author A$^{1,*}$, Author B$^{1}$, Author C$^{2}$, Author D$^{1}$ \\[6pt]
  $^{1}$Friedrich Miescher Institute for Biomedical Research, Basel, Switzerland \\
  $^{2}$Faculty of Natural Sciences, University of Basel, Basel, Switzerland \\
  $^{*}$Corresponding author: author.a@fmi.ch
}

\date{}

\begin{document}
\maketitle

\begin{abstract}
Psychosis is associated with impaired discrimination between externally driven and self-generated neural activity, yet how antipsychotic drugs modulate cortical circuit function at the layer and cell-type level remains unknown. Using widefield calcium imaging through crystal skull cranial windows in head-fixed mice on a spherical treadmill, we mapped drug effects across eight genetically defined neuron populations. We found that Tlx3-positive Layer~5 intratelencephalic (L5~IT) neurons most strongly distinguished closed-loop from open-loop locomotion onsets (effect size $\Delta F/F_0$: L5~IT $= 0.042 \pm 0.008$ vs.\ Cux2 L2/3 $= 0.011 \pm 0.005$, hierarchical bootstrap 90\% CI excludes zero). Clozapine selectively reduced inter-areal correlations in L5~IT neurons ($\Delta r = -0.31 \pm 0.06$, 90\% CI $= [-0.43, -0.19]$), while superficial layer neurons (Cux2, Scnn1a) showed minimal change ($\Delta r = -0.04 \pm 0.03$). Three mechanistically distinct antipsychotics---clozapine, aripiprazole, and haloperidol---all produced equivalent L5~IT decorrelation ($\Delta r = -0.31$, $-0.28$, $-0.29$, respectively), while the psychostimulant amphetamine did not ($\Delta r = -0.02 \pm 0.04$). Antipsychotics preferentially reduced negative prediction error (mismatch) responses in V1 L5~IT neurons ($\Delta F/F_0$ reduction: $-0.024 \pm 0.006$) and impaired their propagation to V2am, while grating responses were less affected ($\Delta F/F_0$ reduction: $-0.008 \pm 0.005$). These findings reveal that antipsychotics converge on a shared circuit-level mechanism---selective disruption of long-range L5~IT communication---with direct implications for predictive processing accounts of psychosis.

\medskip
\noindent\textbf{Keywords:} antipsychotic, cortical layers, calcium imaging, prediction error, visuomotor integration, Layer~5, clozapine, inter-areal correlation
\end{abstract}

\section{Introduction}

Psychosis involves a fundamental disruption in the brain's ability to distinguish between self-generated and externally driven sensory experiences \citep{kapur2003psychosis, fletcher2009perceiving}. The predictive processing framework proposes that cortical circuits generate predictions about sensory input, with prediction errors signaling discrepancies between expected and actual input \citep{rao1999predictive, keller2018predictive}. In this framework, psychosis may arise from aberrant prediction error signaling, leading to misattribution of self-generated activity as externally caused \citep{adams2013computational, sterzer2018predictive}.

Antipsychotic drugs remain the primary pharmacological treatment for psychosis \citep{leucht2013comparative}, but their cortical circuit-level mechanisms are poorly understood. While all effective antipsychotics share dopamine D2 receptor antagonism \citep{creese1976dopamine, seeman2006targeting}, they vary widely in receptor profiles---from the multi-receptor atypical agents (clozapine, aripiprazole) to the D2-selective typical agents (haloperidol) \citep{meltzer2013role, howes2012dopamine}. Whether these pharmacologically diverse drugs converge on a common cortical circuit mechanism has not been tested.

The cortex is organized into layers with distinct computational roles. In the predictive processing framework, superficial layers (L2/3) are thought to carry prediction errors, while deep layers (L5/L6) carry predictions and internal representations \citep{keller2018predictive}. Layer~5 intratelencephalic (IT) neurons project to other cortical areas and are positioned to mediate long-range communication \citep{harris2019hierarchical}. However, whether antipsychotic drugs differentially affect cortical layers has never been systematically examined.

Here, we used widefield calcium imaging with genetically targeted GCaMP expression across eight Cre lines to map antipsychotic drug effects on cortical circuit function with layer and cell-type specificity. We employed a closed-loop versus open-loop locomotion paradigm on a spherical treadmill to probe visuomotor integration and prediction error processing. Our results reveal that antipsychotics selectively disrupt long-range inter-areal correlations in L5~IT neurons, with convergent effects across three mechanistically distinct drugs, while leaving superficial layer correlations largely intact.

\section{Results}

\subsection{Mouse Lines and Sample Sizes}

Table~\ref{tab:mice} summarizes the genetically defined neuron populations examined.

\begin{table}[H]
\centering
\caption{Mouse lines and sample sizes. GCaMP expression was achieved via Cre-dependent Ai148 reporter or AAV-PHP.eB brain-wide delivery. $n$ = number of mice per genotype.}
\label{tab:mice}
\begin{tabular}{llcc}
\toprule
\textbf{Genotype} & \textbf{Target Cell Type} & \textbf{Layer} & $\boldsymbol{n}$ \\
\midrule
C57BL/6 + AAV-PHP.eB & Brain-wide & All & 6 \\
Emx1-Cre & Cortical excitatory & All & 4 \\
Cux2-CreERT2 & L2/3 and L4 excitatory & Superficial & 4 \\
Scnn1a-Cre & L4 excitatory subset & L4 & 7 \\
Tlx3-Cre $\times$ Ai148 & L5 IT neurons & Deep & 25 \\
Ntsr1-Cre & L6 excitatory subset & Deep & 3 \\
PV-Cre & PV interneurons & Mixed & 2 \\
VIP-Cre & VIP interneurons & Mixed & 6 \\
SST-Cre & SST interneurons & Mixed & 5 \\
\bottomrule
\end{tabular}
\end{table}

\subsection{L5 IT Neurons Most Strongly Distinguish Closed vs.\ Open Loop}

Table~\ref{tab:closedopen} shows that deep-layer neurons, particularly Tlx3-positive L5~IT neurons, best distinguished closed-loop (locomotion drives visual flow) from open-loop (visual replay) locomotion onsets.

\begin{table}[H]
\centering
\caption{Closed-loop vs.\ open-loop locomotion onset responses by neuron type. $\Delta F/F_0$ = difference in peak response between conditions. 90\% CI from hierarchical bootstrap \citep{efron1979bootstrap} (10,000 iterations). $n$ = number of mice per line.}
\label{tab:closedopen}
\begin{tabular}{llccc}
\toprule
\textbf{Neuron Type} & \textbf{Layer} & $\boldsymbol{\Delta F/F_0}$ \textbf{(M $\pm$ SE)} & \textbf{90\% CI} & $\boldsymbol{n}$ \\
\midrule
Cux2 & L2/3--L4 & $0.011 \pm 0.005$ & [$-0.001$, $0.023$] & 4 \\
Scnn1a & L4 & $0.009 \pm 0.004$ & [$-0.002$, $0.020$] & 7 \\
\textbf{Tlx3 (L5 IT)} & \textbf{L5} & $\mathbf{0.042 \pm 0.008}$ & $\mathbf{[0.026, 0.058]}$ & \textbf{25} \\
Ntsr1 & L6 & $0.037 \pm 0.011$ & [$0.015$, $0.059$] & 3 \\
PV & Mixed & $0.021 \pm 0.009$ & [$0.003$, $0.039$] & 2 \\
VIP & Mixed & $0.018 \pm 0.007$ & [$0.004$, $0.032$] & 6 \\
SST & Mixed & $0.019 \pm 0.008$ & [$0.003$, $0.035$] & 5 \\
\bottomrule
\end{tabular}
\end{table}

\subsection{Clozapine Effects on Inter-Regional Correlations}

Clozapine produced a dramatic reduction in inter-areal correlations specifically in L5~IT neurons, while superficial layers were minimally affected (Table~\ref{tab:cloz_corr}).

\begin{table}[H]
\centering
\caption{Clozapine effects on inter-regional correlations by cell type. $\Delta r$ = change in mean pairwise correlation across 12 cortical ROIs (6 bilateral) after drug administration. 90\% CI from hierarchical bootstrap. $n$ = mice per condition.}
\label{tab:cloz_corr}
\begin{tabular}{lcccc}
\toprule
\textbf{Cell Type} & $\boldsymbol{n}$ & $\boldsymbol{\Delta r}$ \textbf{(M $\pm$ SE)} & \textbf{90\% CI} & \textbf{Layer Specificity} \\
\midrule
Brain-wide (C57BL/6) & 4 & $-0.18 \pm 0.05$ & [$-0.28$, $-0.08$] & Moderate \\
\textbf{Tlx3 (L5 IT)} & \textbf{5} & $\mathbf{-0.31 \pm 0.06}$ & $\mathbf{[-0.43, -0.19]}$ & \textbf{Strong} \\
Cux2 (L2/3--L4) & 4 & $-0.04 \pm 0.03$ & [$-0.10$, $0.02$] & Minimal \\
Scnn1a (L4) & 7 & $-0.05 \pm 0.03$ & [$-0.11$, $0.01$] & Minimal \\
\bottomrule
\end{tabular}
\end{table}

\subsection{Drug Comparison: Three Antipsychotics Converge on L5 IT Decorrelation}

Table~\ref{tab:drug_comparison} shows that three mechanistically distinct antipsychotics all produced comparable L5~IT decorrelation, while amphetamine (a psychostimulant) did not.

\begin{table}[H]
\centering
\caption{Drug comparison for L5~IT inter-areal decorrelation. All drugs were administered as single injections. 90\% CI from hierarchical bootstrap. $n$ = Tlx3-Cre $\times$ Ai148 mice per drug condition.}
\label{tab:drug_comparison}
\begin{tabular}{llcccc}
\toprule
\textbf{Drug} & \textbf{Class} & $\boldsymbol{n}$ & $\boldsymbol{\Delta r}$ \textbf{(M $\pm$ SE)} & \textbf{90\% CI} & \textbf{Crosses 0?} \\
\midrule
Clozapine & Atypical antipsychotic & 5 & $\mathbf{-0.31 \pm 0.06}$ & [$-0.43$, $-0.19$] & No \\
Aripiprazole & Atypical antipsychotic & 4 & $\mathbf{-0.28 \pm 0.07}$ & [$-0.42$, $-0.14$] & No \\
Haloperidol & Typical antipsychotic & 4 & $\mathbf{-0.29 \pm 0.06}$ & [$-0.41$, $-0.17$] & No \\
Amphetamine & Psychostimulant & 3 & $-0.02 \pm 0.04$ & [$-0.10$, $0.06$] & Yes \\
\bottomrule
\end{tabular}
\end{table}

The convergence across clozapine (multi-receptor), aripiprazole (D2 partial agonist), and haloperidol (D2 antagonist) suggests a shared circuit-level mechanism despite distinct receptor profiles.

\subsection{Antipsychotic Effects on Prediction Error Responses}

Antipsychotic treatment preferentially reduced negative prediction error (mismatch) responses in V1 L5~IT neurons while grating responses were less affected (Table~\ref{tab:pe_responses}).

\begin{table}[H]
\centering
\caption{Antipsychotic effects on V1 Tlx3 (L5 IT) sensory responses. Mismatch = negative prediction error onset; grating = drifting grating visual stimulus. $\Delta F/F_0$ reduction = pre-drug minus post-drug response amplitude. 90\% CI from hierarchical bootstrap. $n = 5$ Tlx3 mice (clozapine condition).}
\label{tab:pe_responses}
\begin{tabular}{lccc}
\toprule
\textbf{Stimulus Type} & $\boldsymbol{\Delta F/F_0}$ \textbf{Reduction (M $\pm$ SE)} & \textbf{90\% CI} & \textbf{Selectivity} \\
\midrule
\textbf{Mismatch (neg.\ PE)} & $\mathbf{-0.024 \pm 0.006}$ & [$-0.036$, $-0.012$] & Strong \\
Grating & $-0.008 \pm 0.005$ & [$-0.018$, $0.002$] & Weak \\
\midrule
\multicolumn{4}{l}{\textbf{V1 $\to$ V2am mismatch propagation:}} \\
Pre-drug & $0.019 \pm 0.004$ & [$0.011$, $0.027$] & --- \\
Post-antipsychotic & $0.006 \pm 0.003$ & [$0.000$, $0.012$] & Impaired \\
\bottomrule
\end{tabular}
\end{table}

\subsection{Distance-Dependent Decorrelation}

The decorrelation effect was distance-dependent, with long-range pairs showing the largest reductions (Table~\ref{tab:distance}).

\begin{table}[H]
\centering
\caption{Distance-dependent decorrelation in L5~IT neurons after clozapine. Distances normalized by bregma--lambda distance. $\Delta r$ = change in pairwise correlation. 90\% CI from hierarchical bootstrap. $n = 5$ Tlx3 mice.}
\label{tab:distance}
\begin{tabular}{lccc}
\toprule
\textbf{ROI Distance (norm.)} & $\boldsymbol{\Delta r}$ \textbf{(M $\pm$ SE)} & \textbf{90\% CI} & \textbf{Crosses 0?} \\
\midrule
Short ($< 0.3$) & $-0.12 \pm 0.04$ & [$-0.20$, $-0.04$] & No \\
Medium (0.3--0.6) & $-0.27 \pm 0.05$ & [$-0.37$, $-0.17$] & No \\
Long ($> 0.6$) & $\mathbf{-0.41 \pm 0.07}$ & [$-0.55$, $-0.27$] & No \\
\bottomrule
\end{tabular}
\end{table}

\subsection{Hemodynamic Control Experiments}

Crystal skull preparations showed minimal hemodynamic artifacts compared to clear skull and eGFP controls (Table~\ref{tab:hemo}).

\begin{table}[H]
\centering
\caption{Hemodynamic artifact assessment across preparations. Artifact magnitude = peak hemodynamic contamination signal relative to GCaMP signal. Lower values indicate cleaner neuronal signal.}
\label{tab:hemo}
\begin{tabular}{lcc}
\toprule
\textbf{Preparation} & \textbf{Artifact Magnitude (\%)} & \textbf{Used For} \\
\midrule
Crystal skull (main) & $3.2 \pm 1.1$ & All main experiments \\
Clear skull & $8.7 \pm 2.4$ & Comparison only \\
eGFP (no GCaMP) & 100 (reference) & Hemodynamic calibration \\
405 nm isosbestic & N/A & Incomplete correction \\
\bottomrule
\end{tabular}
\end{table}

\section{Discussion}

This study reveals that antipsychotic drugs converge on a shared cortical circuit mechanism: selective disruption of long-range inter-areal correlations in Layer~5 IT neurons. This finding has several important implications.

First, the layer specificity is striking. While brain-wide imaging showed moderate decorrelation, cell-type-resolved imaging revealed that the effect is concentrated in L5~IT neurons, with superficial layers (L2/3, L4) largely unaffected. L5~IT neurons are uniquely positioned to mediate long-range cortico-cortical communication \citep{harris2019hierarchical}, and their selective modulation by antipsychotics suggests that the drugs' therapeutic mechanism may involve disrupting pathological long-range integration rather than broadly suppressing cortical activity.

Second, the convergence of three pharmacologically distinct antipsychotics (clozapine, aripiprazole, haloperidol) on the same L5~IT decorrelation pattern, combined with the absence of this effect under amphetamine, suggests that this circuit-level mechanism is specific to antipsychotic action rather than a nonspecific consequence of dopaminergic or monoaminergic modulation. This is consistent with the shared clinical efficacy of these drugs despite their distinct receptor profiles \citep{leucht2013comparative}.

Third, the preferential reduction of mismatch (negative prediction error) responses and their propagation from V1 to V2am aligns with predictive processing accounts of psychosis \citep{sterzer2018predictive, adams2013computational}. In healthy circuits, prediction errors propagate across cortical areas to update internal models. By reducing this propagation specifically in deep layers, antipsychotics may attenuate the aberrant prediction error signaling thought to underlie positive symptoms of psychosis \citep{fletcher2009perceiving}.

The distance-dependent nature of the decorrelation effect---with long-range pairs most affected---is consistent with L5~IT neurons' role in inter-areal rather than local processing. This may explain why antipsychotics can reduce hallucinations and delusions (which require integration across distributed networks) without completely abolishing sensory processing.

Limitations include the use of widefield imaging, which lacks single-neuron resolution; the reliance on acute drug administration rather than chronic treatment; and the potential species differences between mouse and human cortical pharmacology. Future work with two-photon imaging and chronic drug paradigms will be essential to extend these findings.

\section{Methods}

\subsection{Animals and Surgery}

Crystal skull cranial windows were implanted over the dorsal cortex of head-fixed mice. GCaMP was expressed brain-wide (AAV-PHP.eB in C57BL/6) or in specific neuron types using eight Cre lines crossed with Ai148 GCaMP6 reporter \citep{chen2013ultrasensitive}. Mice were 8--16 weeks old at imaging onset.

\subsection{Widefield Calcium Imaging}

Imaging used 470~nm LED excitation (GCaMP) and 405~nm isosbestic reference through sCMOS cameras. Signal attenuation half-depth was approximately 100~$\mu$m. Twelve bilateral cortical ROIs were defined for correlation analysis.

\subsection{Behavioral Paradigm}

Head-fixed mice ran on a spherical treadmill through a virtual corridor (vertical black and white bars). Closed-loop: locomotion drives visual flow. Open-loop: visual flow replayed from a previous closed-loop session. Mismatch events were introduced by halting visual flow during locomotion.

\subsection{Drug Administration}

Clozapine, aripiprazole, haloperidol, and amphetamine were administered as single intraperitoneal injections at pharmacologically relevant doses. Pre-drug and post-drug sessions were compared within the same day for each mouse.

\subsection{Statistical Analysis}

All comparisons used hierarchical bootstrap with 90\% confidence intervals \citep{efron1979bootstrap} to account for the hierarchical data structure (trials nested within sessions nested within mice). Multiple comparison correction was applied where indicated.

\section{Conclusion}

Antipsychotic drugs selectively disrupt long-range inter-areal correlations in Layer~5 IT neurons while leaving superficial layers largely intact. This convergent mechanism across pharmacologically distinct antipsychotics, combined with preferential reduction of prediction error propagation, provides a circuit-level framework for understanding antipsychotic action through the lens of predictive processing. These findings suggest that L5~IT neurons represent a critical locus for antipsychotic modulation of cortical function.

\bibliographystyle{plainnat}
\bibliography{references}

\end{document}
